\section{Residue fields and the discrete valuation topology}

\begin{theorem}[Residue field and completion]
    Let \((K,\varphi)\) be a valued field with valuation ring
    \[
        R=\{x\in K\mid \varphi(x)\leq 1\}
    \]
    and valuation ideal
    \[
        P=\{x\in K\mid \varphi(x)<1\}.
    \]
    Let \((\widehat{K},\widehat{\varphi})\) be the completion of \(K\), and let
    \[
        \widehat{R}
        =
        \{x\in \widehat{K}\mid \widehat{\varphi}(x)\leq 1\}
    \]
    be its valuation ring, with valuation ideal
    \[
        \widehat{P}
        =
        \{x\in \widehat{K}\mid \widehat{\varphi}(x)<1\}.
    \]
    Then there is a canonical isomorphism of residue fields
    \[
        R/P \cong \widehat{R}/\widehat{P}.
    \]
\end{theorem}

\begin{proof}
    We identify \(K\) with its image in \(\widehat{K}\). Since the embedding
    \[
        K\hookrightarrow \widehat{K}
    \]
    preserves the valuation, we have
    \[
        \widehat{\varphi}|_K=\varphi.
    \]
    Hence
    \[
        R\subseteq \widehat{R},
        \qquad
        P=R\cap \widehat{P}.
    \]

    It remains to show that
    \[
        \widehat{R}=R+\widehat{P}.
    \]
    The inclusion
    \[
        R+\widehat{P}\subseteq \widehat{R}
    \]
    is clear, since \(R\subseteq \widehat{R}\) and \(\widehat{P}\subseteq \widehat{R}\).

    Conversely, let \(a\in \widehat{R}\). If \(a=0\), then \(a\in R+\widehat{P}\).
    Suppose \(a\neq 0\). Since \(K\) is dense in \(\widehat{K}\), there exists a sequence
    \[
        (a_n)_{n\geq 1}
    \]
    in \(K\) such that
    \[
        \lim_{n\to\infty} a_n=a.
    \]
    Thus
    \[
        \lim_{n\to\infty}\widehat{\varphi}(a_n-a)=0.
    \]
    Since \(a\neq 0\), we have
    \[
        \widehat{\varphi}(a)>0.
    \]
    Therefore, for some sufficiently large \(n\),
    \[
        \widehat{\varphi}(a_n-a)<\widehat{\varphi}(a).
    \]
    Since
    \[
        a_n=a+(a_n-a)
    \]
    and
    \[
        \widehat{\varphi}(a_n-a)<\widehat{\varphi}(a),
    \]
    the strong triangle inequality for a non-archimedean multiplicative valuation implies
    \[
        \widehat{\varphi}(a_n)
        =
        \widehat{\varphi}\bigl(a+(a_n-a)\bigr)
        =
        \widehat{\varphi}(a).
    \]
    Indeed, for a multiplicative non-archimedean valuation, if
    \[
        \widehat{\varphi}(x)<\widehat{\varphi}(y),
    \]
    then
    \[
        \widehat{\varphi}(x+y)=\widehat{\varphi}(y).
    \]
    Since \(a\in \widehat{R}\), we have
    \[
        \widehat{\varphi}(a)\leq 1.
    \]
    Hence
    \[
        \varphi(a_n)=\widehat{\varphi}(a_n)\leq 1,
    \]
    so
    \[
        a_n\in R.
    \]
    Moreover,
    \[
        \widehat{\varphi}(a-a_n)<\widehat{\varphi}(a)\leq 1,
    \]
    hence
    \[
        a-a_n\in \widehat{P}.
    \]
    Therefore
    \[
        a=a_n+(a-a_n)\in R+\widehat{P}.
    \]
    Thus
    \[
        \widehat{R}\subseteq R+\widehat{P}.
    \]
    Hence
    \[
        \widehat{R}=R+\widehat{P}.
    \]

    Now by the second isomorphism theorem,
    \[
        \widehat{R}/\widehat{P}
        =
        (R+\widehat{P})/\widehat{P}
        \cong
        R/(R\cap \widehat{P}).
    \]
    Since
    \[
        R\cap \widehat{P}=P,
    \]
    we obtain
    \[
        \widehat{R}/\widehat{P}\cong R/P.
    \]
\end{proof}

\begin{proposition}[Fundamental system of neighbourhoods]
    Let \((K,v)\) be a field with a normalized discrete valuation. Let
    \[
        R=\{x\in K\mid v(x)\geq 0\}
    \]
    be its valuation ring, and let
    \[
        P=\{x\in K\mid v(x)>0\}
    \]
    be its valuation ideal. Then
    \[
        P=J(R),
    \]
    and the decreasing sequence of ideals
    \[
        R=P^0\supset P\supset P^2\supset \cdots \supset P^\ell\supset \cdots
    \]
    forms a fundamental system of neighbourhoods of \(0\) in \(K\).
\end{proposition}

\begin{proof}
    Since \(v\) is a normalized discrete valuation, we have
    \[
        v(K^\times)=\mathbb{Z}.
    \]
    Choose a uniformizer \(\pi\in R\), so that
    \[
        v(\pi)=1.
    \]
    Then
    \[
        P=(\pi),
    \]
    and more generally
    \[
        P^\ell=(\pi^\ell)
        =
        \{x\in R\mid v(x)\geq \ell\}
        \qquad
        \text{for every } \ell\geq 0.
    \]
    Thus saying that an element is sufficiently close to \(0\) is equivalent to
    saying that it lies in \(P^\ell\) for large \(\ell\). Hence the ideals
    \[
        P^\ell,\qquad \ell\geq 0,
    \]
    form a fundamental system of neighbourhoods of \(0\).
\end{proof}

\begin{corollary}[Cauchy criterion]
    Let \((a_n)_{n\geq 1}\) be a sequence in \(K\). Then \((a_n)\) is a Cauchy
    sequence if and only if for every natural number \(\ell\geq 1\), there exists
    a natural number \(N\geq 1\) such that
    \[
        a_m-a_n\in P^\ell
        \qquad
        \text{whenever } m,n\geq N.
    \]
\end{corollary}

\begin{proof}
    By definition, \((a_n)\) is Cauchy if and only if for every neighbourhood
    \(U\) of \(0\), there exists \(N\geq 1\) such that
    \[
        a_m-a_n\in U
        \qquad
        \text{for all } m,n\geq N.
    \]
    Since the ideals \(P^\ell\) form a fundamental system of neighbourhoods of
    \(0\), it is enough to test this condition on \(U=P^\ell\). Hence
    \((a_n)\) is Cauchy if and only if for every \(\ell\geq 1\), there exists
    \(N\geq 1\) such that
    \[
        a_m-a_n\in P^\ell
        \qquad
        \text{for all } m,n\geq N.
    \]
\end{proof}

\begin{corollary}[Convergence criterion]
    Let \((a_n)_{n\geq 1}\) be a sequence in \(K\), and let \(a\in K\). Then
    \[
        \lim_{n\to\infty} a_n=a
    \]
    if and only if for every natural number \(\ell\geq 1\), there exists a natural
    number \(N\geq 1\) such that
    \[
        a_n-a\in P^\ell
        \qquad
        \text{whenever } n\geq N.
    \]
\end{corollary}

\begin{proof}
    The sequence \((a_n)\) converges to \(a\) if and only if
    \[
        a_n-a\longrightarrow 0.
    \]
    Since the ideals \(P^\ell\) form a fundamental system of neighbourhoods of
    \(0\), this is equivalent to saying that for every \(\ell\geq 1\), there exists
    \(N\geq 1\) such that
    \[
        a_n-a\in P^\ell
        \qquad
        \text{for all } n\geq N.
    \]
\end{proof}

\begin{lemma}[Equivalent discrete multiplicative valuations]
    Let \(\varphi_1\) and \(\varphi_2\) be discrete multiplicative valuations of a field
    \(K\). The following conditions are equivalent:
    \begin{enumerate}
        \item \(\varphi_1\sim \varphi_2\), that is, there exists \(\lambda>0\) such that
              \[
                  \varphi_2(x)=\varphi_1(x)^\lambda
                  \qquad
                  \text{for all }x\in K.
              \]
        \item The valuation ideals are the same:
              \[
                  \{\,x\in K\mid \varphi_1(x)<1\,\}
                  =
                  \{\,x\in K\mid \varphi_2(x)<1\,\}.
              \]
        \item The topologies on \(K\) induced by \(\varphi_1\) and \(\varphi_2\) are the same.
    \end{enumerate}
\end{lemma}

\begin{proof}
    For \(i=1,2\), let
    \[
        R_i:=\{\,x\in K\mid \varphi_i(x)\leq 1\,\}
    \]
    be the valuation ring of \(\varphi_i\), and let
    \[
        P_i:=\{\,x\in K\mid \varphi_i(x)<1\,\}
    \]
    be its valuation ideal.

    We first note that condition \((2)\) is precisely the assertion
    \[
        P_1=P_2.
    \]

    \((1)\Rightarrow (2)\).
    Suppose that \(\varphi_2=\varphi_1^\lambda\) for some \(\lambda>0\). Then, for every
    \(x\in K\),
    \[
        \varphi_2(x)<1
        \Longleftrightarrow
        \varphi_1(x)^\lambda<1
        \Longleftrightarrow
        \varphi_1(x)<1.
    \]
    Hence \(P_1=P_2\).

    \((2)\Rightarrow (1)\).
    Assume that
    \[
        P_1=P_2=:P.
    \]
    Since \(\varphi_1\) and \(\varphi_2\) are discrete valuations, their valuation rings
    are discrete valuation rings with maximal ideal \(P\). Hence the common ideal \(P\)
    is principal. Choose a uniformizer
    \[
        \pi\in P.
    \]
    Then every element \(x\in K^\times\) admits a unique expression
    \[
        x=\pi^n u,
        \qquad n\in \mathbb Z,\quad u\in R_i^\times.
    \]
    Therefore, for \(i=1,2\),
    \[
        \varphi_i(x)=\varphi_i(\pi)^n.
    \]
    Put
    \[
        \rho_i:=\varphi_i(\pi).
    \]
    Since \(\pi\in P\), we have \(0<\rho_i<1\). Choose
    \[
        \lambda:=\frac{\log \rho_2}{\log \rho_1}>0.
    \]
    Then
    \[
        \rho_2=\rho_1^\lambda.
    \]
    Hence for every \(x=\pi^n u\in K^\times\),
    \[
        \varphi_2(x)
        =
        \rho_2^n
        =
        (\rho_1^\lambda)^n
        =
        (\rho_1^n)^\lambda
        =
        \varphi_1(x)^\lambda.
    \]
    Thus
    \[
        \varphi_2=\varphi_1^\lambda,
    \]
    so \(\varphi_1\sim \varphi_2\).

    \((1)\Rightarrow (3)\).
    If \(\varphi_2=\varphi_1^\lambda\) with \(\lambda>0\), then for any sequence
    \((x_n)\) in \(K\),
    \[
        \varphi_1(x_n)\to 0
        \Longleftrightarrow
        \varphi_1(x_n)^\lambda\to 0
        \Longleftrightarrow
        \varphi_2(x_n)\to 0.
    \]
    Therefore the two valuations define the same convergent sequences to \(0\), and hence
    induce the same topology on \(K\).

    \((3)\Rightarrow (2)\).
    Assume that the topologies induced by \(\varphi_1\) and \(\varphi_2\) are the same.
    Let \(a\in P_1\). Then
    \[
        \varphi_1(a)<1.
    \]
    Hence
    \[
        \lim_{n\to\infty}\varphi_1(a^n)
        =
        \lim_{n\to\infty}\varphi_1(a)^n
        =
        0.
    \]
    Thus
    \[
        a^n\longrightarrow 0
    \]
    with respect to the topology induced by \(\varphi_1\). Since the two topologies are
    the same, we also have
    \[
        a^n\longrightarrow 0
    \]
    with respect to the topology induced by \(\varphi_2\). Therefore
    \[
        \lim_{n\to\infty}\varphi_2(a^n)
        =
        \lim_{n\to\infty}\varphi_2(a)^n
        =
        0.
    \]
    This implies
    \[
        \varphi_2(a)<1,
    \]
    so \(a\in P_2\). Hence
    \[
        P_1\subseteq P_2.
    \]
    By symmetry, we also get
    \[
        P_2\subseteq P_1.
    \]
    Therefore
    \[
        P_1=P_2.
    \]

    We have proved
    \[
        (1)\Rightarrow (2)\Rightarrow (1),
        \qquad
        (1)\Rightarrow (3),
        \qquad
        (3)\Rightarrow (2).
    \]
    Hence the three conditions are equivalent.
\end{proof}
